\section{IoT Systems}

The recorded lecture is provided \myhref{https://www.youtube.com/watch?v=AP_-JETzseo&feature=youtu.be}{here}.

\begin{ejercicio}
    What is the difference between an IoT system and an IoT platform? Use the WWW to find another IoT platform that is not mentioned in the lecture.\\

    They are different concepts. An IoT system consists of all the HW and SW components that are required to implement an IoT application, and developpers tend to reuse the same components across different applications. This is where the concept of IoT platforms comes into play, because an IoT platform is a framework that provides a set of tools and services to facilitate the development, deployment, and management of IoT applications. By using an IoT platform, developers can focus on the application logic rather than dealing with the complexities of the underlying infrastructure. Therefore, an IoT platform is a tool that can be used to build an IoT system.

    An example of an IoT platform that is not mentioned in the lecture is \myhref{https://www.blynk.io/}{Blynk}. It is an horizontal IoT platform that provides a set of tools and services to facilitate the development, deployment, and management of IoT applications. It allows developers to create custom IoT applications using a drag-and-drop interface, and it supports a wide range of hardware platforms and communication protocols.
\end{ejercicio}

\begin{ejercicio}
    Is the Microsoft Azure IoT Suite a horizontal or a vertical IoT platform? Why?

    It is a horizontal IoT platform because it provides a broad set of features and services that can be customized to fit the specific needs of different applications. It is not designed for a specific industry or use case, but rather it can be used for a wide range of IoT applications across different domains.
\end{ejercicio}

\begin{ejercicio}
    Provide an example (i.e., search on the Internet or come up with your own idea) for a scenario in which you would use predictive data processing but not prescriptive data processing. What is the main legal difference between these?

    There are plenty of cases where predictive data processing can be used without prescriptive data processing for several reasons:
    \begin{itemize}
        \item Prescriptive data processing can simply not be made.
        
        For instance, an IoT application designed to predict the weather conditions in a specific area based on historical data and patterns. The application can provide valuable insights and information about the expected weather conditions, but the weather cannot be changed by the application, so prescriptive data processing is not possible in this case.

        \item Prescriptive data processing can be made, but it is not ethical to do it without human supervision and intervention.
        
        For instance, an IoT application designed to predict the likelihood of a patient developing a certain medical condition based on their health data and lifestyle habits. The application may recommend certain lifestyle changes or medical interventions to reduce the risk of developing the condition, but it is not ethical to implement these recommendations without human supervision.

        \item Prescriptive data processing can be made, but a responsible person must be legally accountable for the consequences of the actions taken based on the prescriptive data processing.
        
        For instance, an IoT application designed to predict the future of the stock market based on historical data and patterns. The application may recommend certain investment strategies to maximize returns, but someone should be legally accountable for the consequences of the actions taken based on these recommendations, as the stock market is highly volatile and unpredictable, and there is a risk of financial loss.
    \end{itemize}

    The main legal difference between predictive and prescriptive data processing is that prescriptive data processing actually performs actions or makes decisions, and they can have legal consequences, so there must be a responsible person who is legally accountable for those consequences.
\end{ejercicio}

\begin{ejercicio}
    What are the differences between mesh-based and edge-based IoT systems? Which one is currently supported by platform providers?\\

    Both IoT system architectures are designed to process data closer to where it is generated, but in the edge-based architecture the IoT things are still really simple and resource-constrained, so they just collect data and send it to the gateway, which is responsible for processing the data and making decisions based on it. In contrast, in the mesh-based architecture, the IoT things are more powerful and can communicate directly with each other, so they can process data and make decisions without the need for a central gateway.

    Currently, even though the most supported architecture by platform providers is the cloud-based one, there are some platforms that support edge-based IoT systems, but mesh-based IoT systems are still not widely supported by platform providers due to the complexity and challenges associated with implementing and managing such systems.
\end{ejercicio}